\subsection{完整决策轨迹与有限池风险}
\label{sec:7-budget-trace}

2,304条采样曲线全部单调，二分均选中枚举的最小值；十二次FULL调用也通过两项检查。
将全池逐查询计数按实际样本ID求和，18,432个网格计数全部复现，
支持在这些子集上使用\cref{eq:6-finite-pool}。
该汇总检查不证明任意子集或位置下的逐查询不变性。
分组路径结构保证与距离路径例外仍按\cref{sec:6-3}限定。

$S=128$时，arXiv/512最大损失0.0050、BGE/512为0.0046；$S=32$最大达0.0294。
因此较小平均损失可以伴随少数明显更差的选择。十二个FULL单元的最大损失仍为0.0006。

arXiv/128在$S=32,64,128$的预测越阈率为0.965／0.875／0.620，
实测0.984／0.875／0.625（\cref{fig:7-finite-pool}）。
36项比较最大差为BGE/512在$S=128$的0.144（预测0.519、实测0.375）。
预测以稳定逐查询计数为条件，并用固定预算的留出质量将接纳预算映射到GT损失；
一致性决策本身不使用标签。

\begin{figure}[tbp]
\centering
\includegraphics[width=\linewidth]{supplement/finite-pool.pdf}
\caption[有限池风险预测与实测]{arXiv损失越阈的有限池预测与观测。误差条是在固定池条件下64次独立抽取
子集的一倍估计标准误，不是未来工作负载的不确定性。}
\label{fig:7-finite-pool}
\end{figure}

十轮交替重复将表面的小预算速度反转解释为低预算计时平台。
该观测不要求所有预算和资源阈值上的运行时间单调。

\subsection{直接校准加复用测量}
\label{sec:7-stream-cost}

四个单元以五轮轮换策略顺序、500个留出批次比较固定64／100、采样和FULL校准。
计时含样本收集、校准、预算设置及首批graph设置；
不含构建、预热和扫描校准。查询池重复使用（\cref{tab:7-stream-cost}）。

\par\medskip\noindent\begin{minipage}{\linewidth}
\centering\captionsetup{hypcap=false}
\captionof{table}[500 个复用批次的直接成本]{500个复用批次的直接总毫秒数，五轮平均；采样与FULL包含预算校准，
每批500查询。质量与计时结果一起讨论。}
\label{tab:7-stream-cost}
\footnotesize\setlength{\tabcolsep}{4pt}
\begin{tabular}{@{}lrrrr@{}}
\toprule
数据集/探测数 & 固定 100 & 固定 64 & 规则 $S=128$ & 全池 \\
\midrule
arXiv/128 & 1224 & 1102 & 1291 & 1680 \\
OpenAI-1536/128 & 1491 & 1264 & 852 & 860 \\
OpenAI-3072/128 & 2757 & 2160 & 1769 & 1783 \\
BGE/512 & 5185 & 5018 & 4921 & 5062 \\
\bottomrule
\end{tabular}

\end{minipage}\par\medskip

OpenAI-3072/128的采样策略每轮均选4，相对固定100将总时间从2,757降到1,769\,ms，
约节省36\%，测得recall同为0.9252。FULL选同一预算，含较高校准成本共1,783\,ms。
OpenAI-1536/128的FULL为860／固定100的1,491\,ms，recall同为0.9280；
采样策略平均852\,ms、recall 0.9276。

相对固定100／64，实测平均累计成本首次正节省在OpenAI-1536为4／6批，
OpenAI-3072为4／9批（\cref{fig:7-stream-cost}）；小数稳态拟合属于另一模型估计。
BGE/512以较低时间换取较低平均recall：采样／固定100为4,921／5,185\,ms、
0.9658／0.9670；FULL为5,062\,ms、0.9664。
arXiv的FULL为质量选择200，比未达目标的固定100／64更慢。
校准价值须同时取决于质量和总复用成本。

\begin{figure}[tbp]
\centering
\includegraphics[width=\linewidth]{supplement/stream-cost.pdf}
\caption[校准后复用的直接累计时间]{含预算校准及首批设置的直接平均累计时间，各曲线重复同一固定留出池。
不同策略recall可不同，因此不能仅按时间作质量匹配排名；配置对应\cref{tab:7-stream-cost}。}
\label{fig:7-stream-cost}
\end{figure}
